% 3.6 空间群操作对 Bloch 函数的作用
\section{空间群操作对 Bloch 函数的作用}

现在我们讨论把空间群算符作用于 Bloch 函数所产生的效果。由于本书采用主动操作，1.5 节关于函数空间算符的讨论给出
\begin{equation*}
\{S \mid \mathbf{w}\}\psi(\mathbf{r}) = \psi(S^{-1}(\mathbf{r} - \mathbf{w}))
\tag{3.6.1}
\end{equation*}
其中我们省去了函数空间算符的下标 \textit{a} 和上横线。

首先证明
\begin{equation*}
\exp(\mathrm{i}\mathbf{k} \cdot S\mathbf{r}) = \exp(\mathrm{i}S^{-1}\mathbf{k} \cdot \mathbf{r}).
\tag{3.6.2}
\end{equation*}
这是因为
\begin{equation*}
\mathbf{k} \cdot S\mathbf{r} = \sum_{i}\sum_{j}k_{i}S_{ij}r_{j} = \sum_{i}\sum_{j}S^{\mathrm{T}}_{ji}k_{i}r_{j} = \sum_{i}\sum_{j}S^{-1}_{ji}k_{i}r_{j},
\tag{3.6.3}
\end{equation*}
因为 $S_{ij}$ 是正交矩阵。这意味着由式 (3.6.1) 与 (3.6.2)，
\begin{equation*}
\begin{aligned}
\{S \mid \mathbf{w}\}\exp(\mathrm{i}\mathbf{k} \cdot \mathbf{r}) &= \exp\left\{\mathrm{i}\mathbf{k} \cdot S^{-1}(\mathbf{r} - \mathbf{w})\right\}\\
&= \exp\left\{\mathrm{i}S\mathbf{k} \cdot (\mathbf{r} - \mathbf{w})\right\},
\end{aligned}
\tag{3.6.4}
\end{equation*}
它几乎就是同一个函数，只是以倒格矢 $S\mathbf{k}$ 标记，且中心位于 $\mathbf{r} = \mathbf{w}$。

类似地可以定义
\begin{equation*}
u_{S\mathbf{k}}(\mathbf{r} - \mathbf{w}) = u_{\mathbf{k}}\left\{S^{-1}(\mathbf{r} - \mathbf{w})\right\} = \{S \mid \mathbf{w}\}u_{\mathbf{k}}(\mathbf{r}),
\tag{3.6.5}
\end{equation*}
其中 $u_{\mathbf{k}}(\mathbf{r})$ 满足式 (3.4.8)，即它具有格子的周期性。联立式 (3.6.4) 与 (3.6.5) 并利用 Bloch 函数的定义
\begin{equation*}
\Psi_{\mathbf{k}}(\mathbf{r}) = \exp(\mathrm{i}\mathbf{k} \cdot \mathbf{r})\,u_{\mathbf{k}}(\mathbf{r}),
\tag{3.4.7}
\end{equation*}
得到
\begin{equation*}
\begin{aligned}
\{S \mid \mathbf{w}\}\psi_{\mathbf{k}}(\mathbf{r}) &= \exp\left\{\mathrm{i}S\mathbf{k} \cdot (\mathbf{r} - \mathbf{w})\right\}u_{S\mathbf{k}}(\mathbf{r} - \mathbf{w})\\
&= \psi_{S\mathbf{k}}(\mathbf{r} - \mathbf{w}),
\end{aligned}
\tag{3.6.6}
\end{equation*}
姑且如此记。又由式 (1.5.24) 可得
\begin{equation*}
\{S \mid \mathbf{w}\}\{R \mid \mathbf{v}\}\psi_{\mathbf{k}}(\mathbf{r}) = \psi_{SR\mathbf{k}}(\mathbf{r} - \mathbf{w} - S\mathbf{v}).
\tag{3.6.7}
\end{equation*}
应当注意 $u_{S\mathbf{k}}$ 也满足式 (3.4.8)，因为若 $\mathbf{t}$ 是一个平移，则 $\mathbf{t}' = S^{-1}\mathbf{t}$ 也是，于是
\begin{equation*}
\begin{aligned}
u_{S\mathbf{k}}(\mathbf{r} - \mathbf{w} + \mathbf{t}) &= u_{\mathbf{k}}\left\{S^{-1}(\mathbf{r} - \mathbf{w} + \mathbf{t})\right\}\\
&= u_{\mathbf{k}}\left\{S^{-1}(\mathbf{r} - \mathbf{w}) + \mathbf{t}'\right\}\\
&= u_{\mathbf{k}}\left\{S^{-1}(\mathbf{r} - \mathbf{w})\right\}\quad\text{（由式 (3.4.8)）}\\
&= u_{S\mathbf{k}}(\mathbf{r} - \mathbf{w}).
\end{aligned}
\tag{3.6.8}
\end{equation*}
这意味着式 (3.6.6) 的右端是一个 Bloch 函数。于是，$\{S \mid \mathbf{w}\}$ 作用在以波矢 $\mathbf{k}$ 标记、中心位于 $\mathbf{r} = 0$ 的 Bloch 函数上，把它变换为以波矢 $S\mathbf{k}$ 标记、中心位于 $\mathbf{r} = \mathbf{w}$ 的 Bloch 函数。尽管式 (3.6.4) 与 (3.6.6) 十分相似，二者之间有一个重要差别。若 $S\mathbf{k} = \mathbf{k}$，则 $\{S \mid \mathbf{w}\}\exp(\mathrm{i}\mathbf{k} \cdot \mathbf{r}) = \exp\left\{\mathrm{i}\mathbf{k} \cdot (\mathbf{r} - \mathbf{w})\right\}$，它是完全相同的函数，只是中心位于 $\mathbf{r} = \mathbf{w}$。然而若 $S\mathbf{k} \equiv \mathbf{k}$（等价），则不一定有 $\{S \mid \mathbf{w}\}\psi_{\mathbf{k}}(\mathbf{r}) = \psi_{\mathbf{k}}(\mathbf{r} - \mathbf{w})$；这是因为 $\psi_{\mathbf{k}}(\mathbf{r})$ 上的标记 $\mathbf{k}$ 只是在等价意义下确定的。一个例子有助于说明这一点。设 $\mathbf{k} = \mathbf{0}$，则 $S\mathbf{k} = \mathbf{k}$ 平凡地成立；再取 $\psi_{0}(\mathbf{r}) = 1 + \exp(\mathrm{i}\mathbf{g} \cdot \mathbf{r})$，于是由式 (3.6.5)，
\begin{equation*}
\begin{aligned}
\{S \mid \mathbf{w}\}\psi_{0}(\mathbf{r}) &= \psi_{0}(S^{-1}\mathbf{r} - S^{-1}\mathbf{w})\\
&= 1 + \exp\left\{\mathrm{i}\mathbf{g} \cdot S^{-1}(\mathbf{r} - \mathbf{w})\right\}\\
&= 1 + \exp\left\{\mathrm{i}S\mathbf{g} \cdot (\mathbf{r} - \mathbf{w})\right\}.
\end{aligned}
\tag{3.6.9}
\end{equation*}
但式 (3.6.9) 的右端等于 $\psi_{0}(\mathbf{r} - \mathbf{w})$ 仅当 $S\mathbf{g} = \mathbf{g}$ 时成立。然而，由于 $S\mathbf{g}$ 与 $\mathbf{g}$ 等价，而后者又与零等价，所以与式 (3.6.9) 右端相联系的波矢仍然是零矢量。

\noindent\textbf{定理 3.6.1}\quad 若 $\psi_{\mathbf{k}}(\mathbf{r})$ 是波矢为 $\mathbf{k}$ 的 Bloch 函数，即 $\{E \mid \mathbf{t}\}\psi_{\mathbf{k}}(\mathbf{r}) = \exp(-\mathrm{i}\mathbf{k} \cdot \mathbf{t})\psi_{\mathbf{k}}(\mathbf{r})$，则 $\{S \mid \mathbf{w}\}\psi_{\mathbf{k}}(\mathbf{r})$ 是波矢为 $S\mathbf{k}$ 的 Bloch 函数。

读者应当明白，这并不是通过把 $\{E \mid \mathbf{t}\}$ 作用于式 (3.6.6) 两端并利用已知事实“$\psi_{S\mathbf{k}}(\mathbf{r} - \mathbf{w})$ 是波矢为 $S\mathbf{k}$ 的 Bloch 函数”就能得到的，因为式 (3.6.6) 是数值关系而非函数关系。要在解析上证明该定理，需要经过下面一系列等式：
\begin{equation*}
\begin{aligned}
\{E \mid \mathbf{t}\}\{S \mid \mathbf{w}\}\psi_{\mathbf{k}}(\mathbf{r}) &= \{S \mid \mathbf{w}\}\psi_{\mathbf{k}}(\mathbf{r} - \mathbf{t})\\
&= \psi_{\mathbf{k}}\left\{S^{-1}(\mathbf{r} - \mathbf{t} - \mathbf{w})\right\} = \psi_{S\mathbf{k}}(\mathbf{r} - \mathbf{t} - \mathbf{w})\\
&= \{E \mid \mathbf{t}\}\psi_{S\mathbf{k}}(\mathbf{r} - \mathbf{w}) = \exp(-\mathrm{i}S\mathbf{k} \cdot \mathbf{t})\psi_{S\mathbf{k}}(\mathbf{r} - \mathbf{w})\\
&= \exp(-\mathrm{i}S\mathbf{k} \cdot \mathbf{t})\{S \mid \mathbf{w}\}\psi_{\mathbf{k}}(\mathbf{r}).
\end{aligned}
\tag{3.6.10}
\end{equation*}
当然，也可以从式 (3.6.6) 出发、借助对 Bloch 周期性的几何直观来论证这个定理。作为定理的推论：若 $S\mathbf{k} = \mathbf{k}$，则 $\{S \mid \mathbf{w}\}\psi_{\mathbf{k}}(\mathbf{r})$ 是波矢为 $\mathbf{k}$ 的 Bloch 函数。我们在下一节将用到这个定理。
